\documentclass{article}

\usepackage{hyperref}

\title{\Large \bf ECE 50877 Introduction to Applied 
Cryptography Project Proposals}

\author{
    {\rm Aayuman Ghangurde}\\
    Purdue University
    \and
    {\rm Sai Palaparthi}\\
    Purdue University
    \and
    {\rm Rakib Ahsan}\\
    Purdue University
    \and
    {\rm Zeke Ulrich}\\
    Purdue University
}

\begin{document}

\begin{titlepage} 
\centering 
\vspace*{\fill} 
{\Large\bfseries ECE 50877 \\ Introduction to Applied Cryptography\\[0.5em] 
Project Proposals\par} 
\vspace{2cm} 
{\large Aayuman Ghangurde\\ Purdue University \par} \vspace{0.8cm} 
{\large Sai Palaparthi\\ Purdue University \par} \vspace{0.8cm} 
{\large Rakib Ahsan\\ Purdue University \par} \vspace{0.8cm} 
{\large Zeke Ulrich\\ Purdue University \par} 
\vspace{2.5cm} 

{\large September 11, 2026\par}

\vspace*{\fill} 

\end{titlepage}
\newpage

\section*{Proposal 1}

% This proposal originates with the authors of 
% \href{https://privsec.dev/}{privsec.dev} and is 
% paraphrased from their article on 
% \href{https://privsec.dev/posts/linux/linux-insecurities/}{insecurities in Linux}.

\paragraph{Title:} TPM for Linux Authentication 

\paragraph{Problem Definition:} Operating systems like Android, iOS, macOS, and Windows
use a hardware security module for many functions, such as secure credential storage,
biometric unlock, and verified/measured boot.

On Linux, it is rather the opposite. Out of the box, most 
systems only have at best secure boot enabled (with a shim),
and perhaps some user applications store keys in the TPM (a 
hardware security module available on most modern x86 systems). On some
configured systems, keys are stored in the TPM to unlock an encrypted disk.

While apps on Linux can store credentials on the TPM, this is typically
of limited utility, as they tend not to involve user authentication.
The situation on other operating systems is different. On iOS, apps can
easily store credentials in the system store, which users can unlock securely
using biometrics. For example, on Windows, password managers can securely stash
a long and complex master password (used to first unlock the password database)
into a system store, then allowing users to unlock the password database using
biometric credentials later.

Projects such as credentialsd deal with userspace integration (a generic interface
for providers and consumers of fido2 passkeys). However, they do not deal with
secret material residing on a hardware module.

Thus we aim to answer how to provide hardware-backed credentials and trusted user verification
as reusable system services on Linux, while allowing existing applications and 
authentication infrastructure to consume those services through standardized interfaces.


\paragraph{Proposed Methodology:} 
We propose to build upon the credentialsd architecture by implementing a new credential
provider backed by the Trusted Platform Module. Rather than storing passkeys as userspace secrets,
the provider would direct the TPM to generate keys and sign challenges. The use of these
credentials can be provided through a local pin, and possibly biometrics.

The TPM standard is vast, and thus, we aim to reverse-engineer, or at least determine to some extent,
how Windows uses the TPM for its Windows Hello solution. Given these findings, and research into the TPM standard,
we will interface with the TPM suitable for our purposes.

Further more, credentials require user verification. We will investigate reasonable ways for user credential
entry, and investigate if there are secure methods available on linux as a secondary goal. Since a trusted operating
system is a prerequisite for credential entry, we propose to build the TPM solution on top of verified and measured boot, for which
there is significant prior work. 

We will evaluate this new implementation under a comparison to Windows Hello, and interoperability with existing FIDO
consumers, unauthorized usage of the private keys, integration
with the standard linux logon process, and recovery behavior.

\paragraph{Expected Findings:} We hope to document how other mainstream operating
systems manage user authentication and secrets. We expect that by securing 
keys in a hardware security module, instead of storing them as userspace secrets,
we can improve the security of credentials on Linux.

We expect to demonstrate that TPM can provide the primitives to store non-exportable FIDO2 passkeys
that be used by a credential provider through credentialsd. We expect that the same mechanism can be used
for Linux logon.

We expect to identify the limits of hardware-backed credential protection, and aim to characterize
the limits of trusted user verification on Linux. We expect to identify the necessary components
for a Windows Hello style authentication system, and the areas in which Linux infrastructure is lacking.


% Even though Dr. Ghodsi said she would approve this if we 
% submitted a proposal, I'm unsure if there's enough novel 
% content for a paper. The research space is saturated and 
% Nmap already has version detection. It might yet be an interesting 
% exercise in side-channel analysis and we could see if 
% software/hardware fingerprinting is a practical application
% of machine learning. 

\newpage

\section*{Proposal 2}

This proposal originates with Dr. Zahra Ghodsi of Purdue University. 

\paragraph{Title:} Differential Timing Analysis for Cryptographic Library Fingerprinting

\paragraph{Problem Definition:} The starting step of exploitation is
determining the target's technology stack. An important part of
this process is identifying the cryptographic libraries and
implementations running on the target. Existing tools can often
identify software through explicit version information or protocol
behavior, but such information may be unavailable or intentionally
hidden. This motivates the following research question:
Can differential timing analysis identify the cryptographic library
and version used by a remote target, despite network and system
noise?

\paragraph{Proposed Methodology:} 
\begin{itemize}
    \item We will construct a testbed containing multiple implementations and versions of commonly used cryptographic libraries. 
    \item We will first collect timing measurements while each implementation performs a controlled suite of cryptographic operations under different inputs and configurations. 
    \item We will analyze the resulting timing distributions to identify features that distinguish implementations.
    \item We will then train and evaluate statistical and machine-learning classifiers using these timing features.
\end{itemize}

\paragraph{Expected Findings:} We expect the number of samples 
required to train accurate classifiers to rise substantially with noise. We also expect machine learning classifiers to outperform statistical classifiers, indicating broader applicability of machine learning techniques to side-channel analysis.   


\end{document}